% load package if document contains at least one citation
\usepackage[maxbibnames=99,sorting=none,style=numeric,backend=biber]
{biblatex}
% enable option to print bibliography as a numbered subsection
\defbibheading{bibsubsection}[\bibname]{%
  \subsection{#1}%
}
% \addbibresource{refs.bib}


\begin{document}
\newgeometry{margin=.5cm}
\thispagestyle{empty}
\begin{landscape}

    %% Project-specific Preambles 
    \settoggle{includelogo}{true}

    \newcommand{\thesubject}{Informatik}
    \newcommand{\thetopic}{Arbeitsplan Game}
    \newcommand{\thelogo}{Logos/FreeFlower.pdf}
    \newcommand{\logosize}{.5\textwidth}

    \lhead[ \leftmark ]{Informatik, Python-Game}
    \rhead[Informatik]{\today, Winterthur}

    \begin{longtable}{p{11cm}|p{3cm}|p{3cm}|p{4cm}|p{4cm}}
        \small
        \textbf{Aufgabe / Unteraufgabe}                            & \textbf{Fällig am}\newline(DD/MM/YYYY) & \textbf{Hilfestellungen\newline}$\rightarrow$ Herr Blum   & \textbf{Main Programmer} & \textbf{Co-Programmer(s)} \\\midrule

        \begin{todolist}
            \item \textbf{Formen} für einen Panzer zeichnen
            \begin{todolist}
                \item Rechteck, halber Kreis, Linie kombinieren
                \item Panzer in verschiedenen Farben zeichnen
                \item Grösse und Start-Position bestimmen
            \end{todolist}
        \end{todolist}            &                                        &                                                            &                          &                                                            \\
        \hline
        \begin{todolist}
            \item \textbf{Bewegung} des Panzers implementieren
            \begin{todolist}
                \item Links- und Rechtsbewegung des \textbf{Panzers} mit Tasten (\texttt{A/D} bzw. Pfeil links/rechts)
                \item Rotation der \textbf{Panzer-Kanone} bei Bewegung (\texttt{W/S} bzw. Pfeil hoch/runter)
                \item Sprung-Funktion mit \texttt{Q/P} (mit Schwerkraft)
            \end{todolist}\end{todolist}         &                                        &                                                            &                          &                                      \\
        \hline
        \begin{todolist}
            \item \textbf{Schussmechanik} implementieren
            \begin{todolist}
                \item Schuss mit Leertaste abfeuern (\texttt{SPACE} für Panzer 1, \texttt{ENTER} für Panzer 2)
                \item Kugeln erstellen und zu Liste von Kugeln hinzufügen
                \item Position der Kugeln aktualisieren
                \item Projektil mit Schwerkraft simulieren
                \item Treffererkennung (Panzer trifft Panzer)
            \end{todolist}
        \end{todolist}               &                                        &                                                            &                          &                                                         \\
        \hline
        \begin{todolist}
            \item \textbf{Lebensanzeige} für Panzer erstellen
            \begin{todolist}
                \item Lebenspunkte für jeden Panzer initialisieren
                \item Lebensanzeige zeichnen (Balken oder Zahl)
                \item Lebenspunkte bei Treffer reduzieren
                \item Spielende bei 0 Lebenspunkten
                \item Panzer als Klasse implementieren
            \end{todolist}
        \end{todolist}          &                                        &                                                            &                          &                                                              \\
        \hline
        \begin{todolist}
            \item \textbf{Terrain} erstellen
            \begin{todolist}
                \item Zufälliges Terrain generieren
                \item Terrain zeichnen
                \item Panzer auf Terrain positionieren
            \end{todolist}
        \end{todolist}                           &                                        & \textbf{Zufälliges Terrain?}                               &                          &                                             \\
        \hline
        \begin{todolist}
            \item \textbf{Soundeffekte} und \textbf{Bilder} hinzufügen
            \begin{todolist}
                \item Schuss-Soundeffekt laden und abspielen
                \item Treffer-Soundeffekt laden und abspielen
                \item Hintergrundmusik hinzufügen
                \item Bild bei Explosion anzeigen (Zusatz: langsames Verschwinden)
            \end{todolist}
        \end{todolist} &                                        & \textbf{Langsames Verschwinden?} $\rightarrow$ Herr Blum? &                                                                                                  \\\hline
        $\cdots$                                                   &                                        &                                                            &                          &                           \\
    \end{longtable}
    \newpage
    \fbox{\textbf{Gruppenname}:\hspace{5cm}}\hfill\fbox{\textbf{Gruppen-Mitglieder}:\hspace{14cm}}
    \begin{longtable}{p{11cm}|p{3cm}|p{3cm}|p{4cm}|p{4cm}}
        \small
        \textbf{Aufgabe / Unteraufgabe} & \textbf{Fällig am}\newline(DD/MM/YYYY) & \textbf{Hilfestellungen\newline}$\rightarrow$ Herr Blum & \textbf{Main Programmer} & \textbf{Co-Programmer(s)} \\\midrule
        \begin{todolist}
            \item $\;$
            \begin{todolist}
                \foreach \subtasks in {1,...,5}{
                        \item $\;$
                    }
            \end{todolist}
        \end{todolist}       &                                        &                                                          &                          &                                      \\
        \hline    \begin{todolist}
                      \item $\;$
                      \begin{todolist}
                \foreach \subtasks in {1,...,5}{
                        \item $\;$
                    }
            \end{todolist}
                  \end{todolist}      &                                        &                                                          &                          &                             \\
        \hline    \begin{todolist}
                      \item $\;$
                      \begin{todolist}
                \foreach \subtasks in {1,...,5}{
                        \item $\;$
                    }
            \end{todolist}
                  \end{todolist}      &                                        &                                                          &                          &                             \\
        \hline    \begin{todolist}
                      \item $\;$
                      \begin{todolist}
                \foreach \subtasks in {1,...,5}{
                        \item $\;$
                    }
            \end{todolist}
                  \end{todolist}      &                                        &                                                          &                          &                             \\
        \hline    \begin{todolist}
                      \item $\;$
                      \begin{todolist}
                \foreach \subtasks in {1,...,5}{
                        \item $\;$
                    }
            \end{todolist}
                  \end{todolist}      &                                        &                                                          &                          &                             \\
    \end{longtable}
    \clearpage

    % Mock-Up
    \noindent
    \begin{tabularx}{.95\pageheight}{X|X}
        % --- First row ---
        % First quadrant: two image placeholders
        \begin{minipage}[t][0.48\textheight][t]{\linewidth}
            \textbf{Game-Zeichnungen}:
            \begin{center}
                \fbox{\parbox[c][0.2\textheight][c]{0.475\linewidth}{$\;$}}
                \fbox{\parbox[c][0.2\textheight][c]{0.475\linewidth}{$\;$}}
                \fbox{\parbox[c][0.2\textheight][c]{0.475\linewidth}{$\;$}}
                \fbox{\parbox[c][0.2\textheight][c]{0.475\linewidth}{$\;$}}
            \end{center}
        \end{minipage}
         &
        % Second quadrant
        \begin{minipage}[t][0.48\textheight][t]{\linewidth}
            \textbf{Challenges:}
        \end{minipage}
        \\
        \midrule
        % --- Second row ---
        % Third quadrant
        \begin{minipage}[t][0.48\textheight][t]{\linewidth}
            \textbf{Beschrieb Grundfunktionen + mögliche Erweiterungen:}

        \end{minipage}
         &
        % Fourth quadrant
        \begin{minipage}[t][0.48\textheight][t]{\linewidth}
            \textbf{Zeitplanung:}
        \end{minipage}
    \end{tabularx}

\end{landscape}
\restoregeometry
\end{document}